\section{自由场理论}
场论的限制\begin{enumerate}
    \item 散射理论中相互作用哈密顿量为
        \begin{equation}
        V(t)=\int d^3x\,\mathscr{H}(\mathbf{x},t)
        \end{equation}
        其中 $\mathscr{H}$ 为洛伦兹标量：
            \begin{equation}
                    U_0(\Lambda,a)\,\mathscr{H}(x)\,U_0^\dagger(\Lambda,a)
                    =\mathscr{H}(\Lambda x + a)
            \end{equation}

        微因果性要求
            \begin{equation}
                [\mathscr{H}(\mathbf{x},t),\mathscr{H}(\mathbf{x}',t)]=0
                \quad\text{当}\quad(x-x')^2 \ge 0
            \end{equation}
    \item 集团分解原理表明：若用产生算子 $a(\boldsymbol{p},\sigma,n)$与湮灭算子$a^\dagger(\boldsymbol{p},\sigma,n)$构造$S$矩阵的连接部分$S^\mathbb{C}\thicksim \mathscr{H}$，则对应的位置空间产生场与湮灭场必须为
            \begin{subequations}
                   \begin{align}
                     &\Psi^{+}_\ell(x)=\sum_{\sigma,n}\int d^3p\;u_\ell(x;\boldsymbol{p},\sigma,n)\,a^\dagger(\boldsymbol{p},\sigma,n)\\
                    &\Psi^{-}_\ell(x)=\sum_{\sigma,n}\int d^3p\;v_\ell(x;\boldsymbol{p},\sigma,n)\,a(\boldsymbol{p},\sigma,n)
                   \end{align}
            \end{subequations}
            对系数$u_\ell(x;\boldsymbol{p},\sigma,n)$和$v_\ell(x;\boldsymbol{p},\sigma,n)$进行选择，可以让$\Psi^\pm$的$Lorentz$变换\footnote{原则上，对于不同的场，有不同的变换矩阵 $D^{\pm}$，我们看到 $D$-矩阵构成了齐次 Lorentz 群的一个表示：
$$ \mathscr{D}(\Lambda_1)\mathscr{D}(\Lambda_2) = D(\Lambda_1\Lambda_2) .  $$
可见 $\mathscr{D}$-矩阵构成了齐次 Lorentz 群的一个表示：洛伦兹群拥有多种维度的不可约表示，包括平凡、旋量、矢量以及更高维的张量表示等。一般而言，处理的 $\mathscr{D}(\Lambda)$ 并非必须是不可约的，它通常呈现为由上述不可约表示组成的分块对角矩阵。所以指标 $\ell$ 在这里是一个复合指标：它既遍历不同的分块（代表不同的粒子种类），也遍历分块内部的具体分量（比如狄拉克旋量和外尔旋量之间的关系）。}有如下形式：
        \begin{subequations}
                   \begin{align}
                     U_0(\Lambda,a)&\Psi^{+}_\ell(x)U_0^\dagger(\Lambda,a)=\sum_{\bar{\ell}}\mathscr{D}_{\ell\bar{\ell}}(\Lambda^{-1})\Psi_{\bar{\ell}}^+(\Lambda x+a)\\
                    U_0(\Lambda,a)&\Psi^{-}_\ell(x)U_0^\dagger(\Lambda,a)=\sum_{\bar{\ell}}\mathscr{D}_{\ell\bar{\ell}}(\Lambda^{-1})\Psi_{\bar{\ell}}^-(\Lambda x+a)
                   \end{align}
            \end{subequations}
    因此相互作用密度可写为:
        \begin{equation}
                \begin{aligned}
                    \mathscr{H}(x)
                    &=
                    \sum_{n,m}
                    \sum_{l'_1\ldots l'_n}
                    \sum_{l_1\ldots l_m}
                    g_{l'_1\ldots l'_n;\,l_1\ldots l_m}
                    \Psi^+_{l_1}\cdots
                    \Psi^+_{l_n}
                    \,
                    \Psi^-_{l'_1}\cdots
                    \Psi^-_{l'_m}
                \end{aligned}
        \end{equation}
    其中系数必须满足
            $g_{l'_1\ldots l'_n;\,l_1\ldots l_m}
            \ \text{为 Lorentz 不变量}$
    从而保证$ \mathscr{H}(x)$为标量.
\end{enumerate}
\color{blue}这里，下面系数$u_\ell(x;\boldsymbol{p},\sigma,n)$和$v_\ell(x;\boldsymbol{p},\sigma,n)$的选择进行讨论\normalcolor 
,产生和湮灭算符的变换如同单粒子态变换一致
\begin{subequations}
    \begin{align}
                &U_0(\Lambda,a)a^\dagger(\boldsymbol{p},\sigma,n)U_0^\dagger(\Lambda,a)=e^{-i(\Lambda p)^\mu a_\mu}\sqrt{\frac{(\Lambda p)^0}{p^0}}\sum_{\bar{\sigma}}D_{\sigma \bar{\sigma}}^{j_n}\left\{W^{-1}(\Lambda,p)\right\}a^\dagger(\Lambda\boldsymbol{p},\bar{\sigma},n)\\
                &U_0(\Lambda,a)a(\boldsymbol{p},\sigma,n)U_0^\dagger(\Lambda,a)=e^{i(\Lambda p)^\mu a_\mu}\sqrt{\frac{(\Lambda p)^0}{p^0}}\sum_{\bar{\sigma}}D_{\sigma \bar{\sigma}}^{j_n*}\left\{W^{-1}(\Lambda,p)\right\}a(\Lambda\boldsymbol{p},\bar{\sigma},n)
    \end{align}
\end{subequations}
所以对于场的$Lorentz$变换：即可以利用产生(湮灭)算符的\text{Poincar\'e}变换进行表示，也可以利用产生(湮灭)场的$Lorentz$变换进行表示，简写就是：
\begin{equation}
    U_0(\Lambda,a)\Psi^{\pm}_\ell(x)U_0^\dagger(\Lambda,a)
    =\left\{
        \begin{aligned}
            &\varpropto U_0(\Lambda,a)\left[\begin{aligned}
                &a^\dagger(\boldsymbol{p},\sigma,n)\\&a(\boldsymbol{p},\sigma,n)
            \end{aligned}\right] U_0^\dagger(\Lambda,a)\\
            &\varpropto \mathscr{D}_{\ell\bar{\ell}}(\Lambda^{-1})\left[\begin{aligned}
                &\Psi_{\bar{\ell}}^+(\Lambda x+a)\\&\Psi_{\bar{\ell}}^-(\Lambda x+a)
            \end{aligned}\right]
        \end{aligned}
    \right.
\end{equation}
利用这两种表述方式，\color{blue}链接场和粒子\normalcolor;可以选择出系数矩阵$u_\ell(x;\boldsymbol{p},\sigma,n)$和$v_\ell(x;\boldsymbol{p},\sigma,n)$的表达形式；\color{blue}事实上，系数函数和两个群的结构息息相关\normalcolor；下面给出决定系数矩阵的表达式的推导：
\vspace{2ex}
\hrule
\vspace{2ex}
\centerline{\text{产生场的系数}}
\begin{align*}
    L.H.S=&\;U_0(\Lambda,a)\Psi^{+}_\ell(x)U_0^\dagger(\Lambda,a)\\
    =&\sum_{\sigma,n}\int d^3p\, u_{\ell}(x;\boldsymbol{p},\sigma,n) U_0(\Lambda,a)a^\dagger(\boldsymbol{p},\sigma,n)U_0^\dagger(\Lambda,a)\\
    =&\sum_{\sigma,n}\int d^3p\, u_{\ell}(x;\boldsymbol{p},\sigma,n)\,e^{-i(\Lambda p)^\mu a_\mu}\sqrt{\frac{(\Lambda p)^0}{p^0}}\sum_{\bar{\sigma}}D_{\sigma \bar{\sigma}}^{j_n}\left\{W^{-1}(\Lambda,p)\right\}a^\dagger(\Lambda\boldsymbol{p},\bar{\sigma},n)\\
    &\xrightarrow{d^3p\to \frac{d^3(\Lambda p)}{(\Lambda p)^0}p^0}\sum_{\sigma\bar{\sigma},n}\int d^3\Lambda p\, u_{\ell}(x;\boldsymbol{p},\sigma,n)\,e^{-i(\Lambda p)^\mu a_\mu}\sqrt{\frac{p^0}{(\Lambda p)^0}}\,D_{\sigma \bar{\sigma}}^{j_n}\left\{W^{-1}(\Lambda,p)\right\} a^\dagger(\Lambda\boldsymbol{p},\bar{\sigma},n)
\end{align*}
\begin{align*}
    R.H.S=&\sum_{\bar{\ell}}\mathscr{D}_{\ell\bar{\ell}}(\Lambda^{-1})\Psi_{\bar{\ell}}^+(\Lambda x+a)\\
    =&\sum_{\bar{\ell}}\sum_{\sigma,n}\int d^3\Lambda p\,\mathscr{D}_{\ell\bar{\ell}}(\Lambda^{-1}) u_{\bar{\ell}}(\Lambda x+a;\Lambda \boldsymbol{p},\sigma,n) a^\dagger(\Lambda \boldsymbol{p},\sigma,n)
\end{align*}
$L.H.S=R.H.S\implies$ 
\begin{equation*}
    \begin{aligned}
    &\sum_{\sigma\bar{\sigma}}u_{\ell}(x;\boldsymbol{p},\sigma,n)\,e^{-i(\Lambda p)^\mu a_\mu}\sqrt{\frac{p^0}{(\Lambda p)^0}}
    D_{\sigma \bar{\sigma}}^{j_n}\left\{W^{-1}(\Lambda,p)\right\}a^\dagger(\Lambda\boldsymbol{p},\bar{\sigma},n)\\[6pt]
    =&\sum_{\bar{\ell}}\sum_{\sigma}\mathscr{D}_{\ell\bar{\ell}}(\Lambda^{-1})u_{\bar{\ell}}(\Lambda x+a;\Lambda \boldsymbol{p},\sigma,n) a^\dagger(\Lambda \boldsymbol{p},\sigma,n)
\end{aligned}
\end{equation*}
\label{u_1}
左式令 $\bar{\sigma}\to\sigma$，$\sigma\to\bar{\sigma}$：
\begin{align*}
    &\sum_{\sigma\bar{\sigma}}u_{\ell}(x;\boldsymbol{p},\bar{\sigma},n)\,e^{-i(\Lambda p)^\mu a_\mu}\sqrt{\frac{p^0}{(\Lambda p)^0}} D_{\sigma \bar{\sigma}}^{j_n}\left\{W^{-1}(\Lambda,p)\right\}a^\dagger(\Lambda\boldsymbol{p},\sigma,n)
    \\
    =&\sum_{\bar{\ell}}\sum_{\sigma}\mathscr{D}_{\ell\bar{\ell}}(\Lambda^{-1})u_{\bar{\ell}}(\Lambda x+a;\Lambda \boldsymbol{p},\sigma,n)a^\dagger(\Lambda \boldsymbol{p},\sigma,n)
\end{align*}
\centerline{\text{湮灭场的系数}}
\begin{align*}
    L.H.S&=\;U_0(\Lambda,a)\Psi^{-}_\ell(x)U_0^\dagger(\Lambda,a)\\
    =&\sum_{\sigma,n}\int d^3p\, v_{\ell}(x;\boldsymbol{p},\sigma,n) U_0(\Lambda,a)a(\boldsymbol{p},\sigma,n)U_0^\dagger(\Lambda,a)\\
    =&\sum_{\sigma,n}\int d^3p\, v_{\ell}(x;\boldsymbol{p},\sigma,n)\,e^{i(\Lambda p)^\mu a_\mu}\sqrt{\frac{(\Lambda p)^0}{p^0}}\sum_{\bar{\sigma}}D_{\sigma \bar{\sigma}}^{j_n*}\left\{W^{-1}(\Lambda,p)\right\}a(\Lambda\boldsymbol{p},\bar{\sigma},n)\\
    &\xrightarrow{d^3p\to \frac{d^3(\Lambda p)}{(\Lambda p)^0}p^0}\sum_{\sigma\bar{\sigma},n}\int d^3\Lambda p\, v_{\ell}(x;\boldsymbol{p},\sigma,n)\,e^{i(\Lambda p)^\mu a_\mu}\sqrt{\frac{p^0}{(\Lambda p)^0}}\,D_{\sigma \bar{\sigma}}^{j_n*}\left\{W^{-1}(\Lambda,p)\right\} a(\Lambda\boldsymbol{p},\bar{\sigma},n)
\end{align*}
\begin{align*}
    R.H.S=&\sum_{\bar{\ell}}\mathscr{D}_{\ell\bar{\ell}}(\Lambda^{-1})\Psi_{\bar{\ell}}^-(\Lambda x+a)\\
    =&\sum_{\bar{\ell}}\sum_{\sigma,n}\int d^3\Lambda p\,\mathscr{D}_{\ell\bar{\ell}}(\Lambda^{-1}) v_{\bar{\ell}}(\Lambda x+a;\Lambda \boldsymbol{p},\sigma,n) a(\Lambda \boldsymbol{p},\sigma,n)
\end{align*}
$L.H.S=R.H.S\implies$ 
\begin{align*}
    &\sum_{\sigma\bar{\sigma}}v_{\ell}(x;\boldsymbol{p},\sigma,n)\,e^{i(\Lambda p)^\mu a_\mu}\sqrt{\frac{p^0}{(\Lambda p)^0}} D_{\sigma \bar{\sigma}}^{j_n*}\left\{W^{-1}(\Lambda,p)\right\}a(\Lambda\boldsymbol{p},\bar{\sigma},n)\\[6pt]
    =&\sum_{\bar{\ell}}\sum_{\sigma}\mathscr{D}_{\ell\bar{\ell}}(\Lambda^{-1})v_{\bar{\ell}}(\Lambda x+a;\Lambda \boldsymbol{p},\sigma,n) a(\Lambda \boldsymbol{p},\sigma,n)
\end{align*}
左式令 $\bar{\sigma}\to\sigma$，$\sigma\to\bar{\sigma}$：
\begin{align*}
    &\sum_{\sigma\bar{\sigma}}v_{\ell}(x;\boldsymbol{p},\bar{\sigma},n)\,e^{i(\Lambda p)^\mu a_\mu}\sqrt{\frac{p^0}{(\Lambda p)^0}} D_{\sigma \bar{\sigma}}^{j_n*}\left\{W^{-1}(\Lambda,p)\right\}a(\Lambda\boldsymbol{p},\sigma,n)
    \\
    =&\sum_{\bar{\ell}}\sum_{\sigma}\mathscr{D}_{\ell\bar{\ell}}(\Lambda^{-1})v_{\bar{\ell}}(\Lambda x+a;\Lambda \boldsymbol{p},\sigma,n) a(\Lambda \boldsymbol{p},\sigma,n)
\end{align*}
最终得到
\begin{equation}
    \begin{aligned}
            \sum_{\bar{\sigma}}
            u_{\ell}(x;\boldsymbol{p},\bar{\sigma},n)
            e^{-i(\Lambda p)^\mu a_\mu}
            \sqrt{\frac{p^0}{(\Lambda p)^0}} 
            &D_{\sigma \bar{\sigma}}^{j_n*}
            \!\left\{W^{-1}(\Lambda,p)\right\} \\
            &=
            \sum_{\bar{\ell}}
            \mathscr{D}_{\ell\bar{\ell}}(\Lambda^{-1})
            u_{\bar{\ell}}(\Lambda x+a;\Lambda \boldsymbol{p},\sigma,n)\\
            %
            \sum_{\bar{\sigma}}
            v_{\ell}(x;\boldsymbol{p},\bar{\sigma},n)
            e^{i(\Lambda p)^\mu a_\mu}
            \sqrt{\frac{p^0}{(\Lambda p)^0}} 
            &D_{\sigma \bar{\sigma}}^{j_n}
            \!\left\{W^{-1}(\Lambda,p)\right\} \\&=
            \sum_{\bar{\ell}}
            \mathscr{D}_{\ell\bar{\ell}}(\Lambda^{-1})
            v_{\bar{\ell}}(\Lambda x+a;\Lambda \boldsymbol{p},\sigma,n)
            \end{aligned}
\end{equation}
\vspace{2ex}
\hrule
\vspace{2ex}
\color{blue}到现在为止，一直在回避一个问题：构建局域场算符时：“粒子”对应 $\text{Poincar\'e}$ 群，而“场”对应$Lorentz$群
\begin{itemize}
    \item \textbf{粒子（物理态）：} 存在于动量空间，必须依赖时空平移来定义其质量(无质量的也是依赖的)和动量属性。
    \item \textbf{场（局域算符）：} 带有内禀指标 $\ell$，定义在时空单点上。平移只是改变了坐标自变量，绝对不会把场的一个分量变成另一个分量。
\end{itemize}


展开系数 $u_\ell$ 和 $v_\ell$ 链接了场和粒子；也就是：粒子态所遵循的 $\text{Poincar\'e}$ 变换强行映射到了场算符所要求的 Lorentz 变换上。\normalcolor 为了清晰地展现这种映射关系，并印证上述二者的内在区别，首先单独考察 $\text{Poincar\'e}$变换中的\textbf{平移变换}：
令$\Lambda=1$,$a=\epsilon$，可以得到
\[\begin{aligned}
    \Lambda=1&\to \mathscr{D}_{\ell\bar{\ell}}(\Lambda^{-1})=1\\
    W(\Lambda,p)=1&\to D_{\sigma \bar{\sigma}}^{j_n}\!\left\{W^{-1}(\Lambda,p)\right\}=1
\end{aligned}\]
所以系数的变换关系变为

\[\begin{aligned}
    u_\ell(x;\boldsymbol{p},\sigma,n)e^{-ip^\mu a_\mu}&=u_{\bar{\ell}}(x+a;\boldsymbol{p},\sigma,n)\\
    &=u_{\bar{\ell}}(\boldsymbol{p},\sigma,n)e^{-i p^\mu (x+a)_\mu}\\
    v_\ell(x;\boldsymbol{p},\sigma,n)e^{ip^\mu a_\mu}&=v_{\bar{\ell}}(x+a;\boldsymbol{p},\sigma,n)\\
    &=v_{\bar{\ell}}(\boldsymbol{p},\sigma,n)e^{i p^\mu (x+a)_\mu}\\
\end{aligned}\]
所以
\[ \begin{aligned}
    &(\frac{1}{2\pi})^{\frac{3}{2}}u_\ell(\boldsymbol{p},\sigma,n)e^{-ip^\mu a_\mu}=(\frac{1}{2\pi})^{\frac{3}{2}}u_{\bar{\ell}}(\boldsymbol{p},\sigma,n)e^{-i p^\mu x_\mu}\\ 
    &(\frac{1}{2\pi})^{\frac{3}{2}}v_\ell(\boldsymbol{p},\sigma,n)e^{ip^\mu a_\mu}=(\frac{1}{2\pi})^{\frac{3}{2}}v_{\bar{\ell}}(\boldsymbol{p},\sigma,n)e^{i p^\mu x_\mu}
\end{aligned}\]
可以看出所谓的场，简单理解就是产生湮灭算符的位置空间形式
\begin{subequations}
    \begin{align}
\Psi_\ell^+(x) &= \sum_{\sigma,n} (2\pi)^{-3/2} \int d^3p \, u_\ell(\boldsymbol{p}, \sigma, n) e^{-ip\cdot x} a^\dagger(\boldsymbol{p}, \sigma, n) , \\
\Psi_\ell^-(x) &= \sum_{\sigma,n} (2\pi)^{-3/2} \int d^3p \, v_\ell(\boldsymbol{p}, \sigma, n) e^{ip\cdot x} a(\boldsymbol{p}, \sigma, n) .
\end{align}
\end{subequations}
我们得出结论：平移在粒子态的表述上是非$trivial$的，但在场的表述上是$trivial$的，将剩下的系数矩阵的变换关系写出，就是：
\begin{subequations}
    \begin{align}
            &\sum_{\bar{\sigma}}
            u_{\ell}(\boldsymbol{p},\bar{\sigma},n)
            \sqrt{\frac{p^0}{(\Lambda p)^0}} 
            D_{\sigma \bar{\sigma}}^{j_n}
            \!\left\{W^{-1}(\Lambda,p)\right\} =
            \sum_{\bar{\ell}}
            \mathscr{D}_{\ell\bar{\ell}}(\Lambda^{-1})
            u_{\bar{\ell}}(\Lambda \boldsymbol{p},\sigma,n)
            \\[10pt]
            %
            &\sum_{\bar{\sigma}}
            v_{\ell}(\boldsymbol{p},\bar{\sigma},n)
            \sqrt{\frac{p^0}{(\Lambda p)^0}} 
            D_{\sigma \bar{\sigma}}^{j_n*}
            \!\left\{W^{-1}(\Lambda,p)\right\} =
            \sum_{\bar{\ell}}
            \mathscr{D}_{\ell\bar{\ell}}(\Lambda^{-1})
            v_{\bar{\ell}}(\Lambda \boldsymbol{p},\sigma,n)
            \end{align}
\end{subequations}
正常接下来应该讨论$Boost$和$Rotation$；但这里先停止，将这两种变换与系数矩阵的关系放在构造场之前.
\vspace{2ex}
\begin{center}
    验证集团分解
\end{center}
\vspace{2ex}
回到集团分解中来，我们得到了场的一般形式，将产生场和湮灭场带入到之前构造的相互作用，并把位置积掉;
\begin{equation}
\begin{aligned}
    V(t)&=\int d^3x\mathscr{H}(x)\sum_{NM} \sum_{\sigma'_1 \cdots \sigma'_N} \sum_{\sigma_1 \cdots \sigma_M} \sum_{n'_1 \cdots n'_N} \sum_{n_1 \cdots n_M} \sum_{\ell'_1 \cdots \ell'_N} \sum_{\ell_1 \cdots \ell_M}\times(2\pi)^{3 - \frac{3N}{2} - \frac{3M}{2}} \\
    &\times\int d^3\mathbf{p}'_1 \cdots d^3\mathbf{p}'_N d^3\mathbf{p}_1 \cdots d^3\mathbf{p}_M  \delta^3(\mathbf{p}'_1 + \cdots  - \mathbf{p}_1 - \cdots- ) \\
    &\times a^\dagger(\mathbf{p}'_1 \sigma'_1 n'_1) \cdots a^\dagger(\mathbf{p}'_N \sigma'_N n'_N) a(\mathbf{p}_M \sigma_M n_M) \cdots a(\mathbf{p}_1 \sigma_1 n_1) \\
    &g_{\ell'_1 \cdots \ell'_N, \, \ell_1 \cdots \ell_M}  u_{\ell'_1}(\mathbf{p}'_1 \sigma'_1 n'_1) \cdots u_{\ell'_N}(\mathbf{p}'_N \sigma'_N n'_N) v_{\ell_1}(\mathbf{p}_1 \sigma_1 n_1) \cdots v_{\ell_M}(\mathbf{p}_M \sigma_M n_M) \ .
\end{aligned}
\end{equation}
\color{blue}可以看出，上式中包含且仅包含一个“唯一”的全局动量守恒 $\delta$ 函数（即 $\delta^3(\mathbf{p}'_1 + \cdots - \mathbf{p}_1 - \cdots)$），而剩下的动量空间系数（由 $g$ 张量和 $u, v$ 展开系数构成）则是充分光滑的函数（至多在零动量处有分支点奇异性，但不包含任何额外的 $\delta$ 函数奇点）。这种“单一 $\delta$ 函数 $\times$ 光滑函数”的漂亮的结构并不是巧合，正是 $S$-矩阵能够满足集团分解原理的的推论事实上，Weinberg 指出我们可以\textbf{反转}这一逻辑来思考：
\begin{itemize}
    \item 任何物理算符都可以写成类似上式中产生/湮灭算符乘积的积分形式；
    \item 集团分解原理强制要求其系数矩阵必须能够提取出\textbf{单个}动量守恒 $\delta$ 函数，剩下一个光滑函数；
    \item 而任何充分光滑的函数，在满足 \textbf{洛伦兹不变性} 的要求下，都可以像上式那样用有限维表示的展开系数 $u, v$ 结合常数张量 $g$ 来表达。
\end{itemize}
\normalcolor
这正是量子场论的核心，\textbf{集团分解}与\textbf{$Lorentz$不变性}，共同迫使相互作用密由局域的产生场和湮灭场来构建。
\vspace{0.5cm}
\begin{center}
    $Lorentz$不变性的要求：微观因果性
\end{center}
\vspace{0.5cm}
微观因果性是：$\text{当}(x-x')^2 \ge 0$有$[\mathscr{H}(\mathbf{x},t),\mathscr{H}(\mathbf{x}',t)]=0$，使用场算符的形式，那么对易关\footnote{$\mp$代表对易关系和反对易}系变为：
\begin{subequations}
\begin{align}
    % 第一行：产生场之间的对易/反对易关系
    &[\Psi^+_\ell(x), \Psi^+_{\bar{\ell}}(x')]_{\mp}= 0 \\
    % 第二行：湮灭场之间的对易/反对易关系
    &[\Psi^-_\ell(x), \Psi^-_{\bar{\ell}}(x')]_{\mp}  = 0 \\
    % 第三行：交叉部分（湮灭场与产生场）
    &[\Psi^-_\ell(x), \Psi^+_{\bar{\ell}}(x')]_{\mp} = 
    \frac{1}{(2\pi)^3}\sum_{\sigma,n}\int d^3 \boldsymbol{p}\,u_\ell(\boldsymbol{p},\sigma,n)v_{\bar{\ell}}(i\boldsymbol{p},\sigma,n)e^{p_\mu(x-x')^\mu}
\end{align}
\end{subequations}
这违背了微观因果性，唯一的解决办法是：将产生场和湮灭场线性组合：即：
\begin{equation}
    \Psi_\ell(x)=\xi_\ell\Psi^+_\ell(x)+\lambda _\ell\Psi^{-}_\ell(x)
\end{equation}
通过调整系数，使得满足对易关系
\begin{equation}
    [\Psi_\ell(x), \Psi_{\bar{\ell}}(x')]_{\mp} = 0
\end{equation}
在之后讨论各类场的时候，再决定系数如何选择；这个式子也叫做因果律。
\begin{center}
    \vspace{2ex}
   \color{black}\text{***}
   \vspace{2ex}
\end{center}
在前面的推导中，我们通过将产生场和湮灭场进行线性组合，构造了满足微观因果性的场 $\Psi(x)$，即在类空间隔 $(x-y)^2 \ge 0$ 下 $[\Psi(x), \Psi(y)] = 0$。然而，为了保证理论满足集团分解原理，相互作用密度必须写成正规乘积（Normal Ordered）的形式：
\begin{equation}
    V = \int d^3x \, N\!\left( \mathscr{F}(\Psi(x), \Psi^\dagger(x)) \right)
\end{equation}
正规乘积操作 $N(\cdots)$ 会强行将所有的产生算符重排到湮灭算符的左边。一个自然的问题是：这种强行重排的操作，是否会破坏构造的场算符对易关系，从而导致微观因果性的失效？计算一个模型：以最简单的双场乘积为例，普通乘积与正规乘积之间仅相差一个缩并项（真空期望值），而这个缩并项纯粹是一个 c-数：
\begin{equation*}
    \Psi(x)\Psi^\dagger(y) = N\!\left( \Psi(x)\Psi^\dagger(y) \right) + \text{c-数}
\end{equation*}
移项可知，正规乘积等价于普通乘积减去一个 c-数。将这一结论推广到任意多项式 $\mathscr{F}$，我们可以将其写为普通乘积与若干 c-数的组合：
\begin{equation*}
    N(\mathscr{F}(x)) = \mathscr{F}(x) - c_1
\end{equation*}
在计算两个时空点处正规乘积算符的对易子时，由于 c-数与任何算符的对易子皆为零，这些 c-数项会自动抵消：
\begin{equation*}
\begin{aligned}
    \left[ N(\mathscr{F}(x)) \,,\, N(\mathscr{F}(y)) \right] 
    &= \left[ \mathscr{F}(x) - c_1 \,,\, \mathscr{F}(y) - c_2 \right] \\
    &= \left[ \mathscr{F}(x) \,,\, \mathscr{F}(y) \right]
\end{aligned}
\end{equation*}
两个正规乘积算符的对易子完全等于它们普通乘积的对易子。既然普通的场多项式 $\mathscr{F}$ 在类空间隔下是对易的，那么它的正规乘积 $N(\mathscr{F})$ 也必然是对易的（要求：相互作用中任意费米场分量的数量为偶数）。 因此正规乘积（保证集团分解原理）与场的线性组合（保证微观因果性）在数学结构上完美自洽。

\vspace{2ex}
\begin{center}
    \large 反粒子
\end{center}
\vspace{2ex}
要满足微观因果性，现在出现了一个阻碍，就是荷守恒：
\begin{itemize}
    \item \textbf{因果性要求}：为了保证微观因果性（即在类空间隔 $[\Psi(x), \Psi^\dagger(y)]_{\pm} = 0$） $\Psi(x)$， 必须同时包含$a$和$a^\dagger$：
    \begin{equation}
        \Psi(x) \sim \int d^3p \, \left( a(\mathbf{p}) e^{-ipx} + a^\dagger(\mathbf{p}) e^{ipx} \right)
    \end{equation}
    \item \textbf{守恒律要求}：S 矩阵的对称性告诉我们理论是存在荷守恒的，所以$S$相互作用哈密顿量 $\mathscr{H}(x)$ 必须由具有确定荷变换性质的场构成：
    \begin{equation}
        [Q, \Psi(x)] = -q \Psi(x)
    \end{equation}
    这意味着场 $\Psi(x)$ 的\textbf{每一项}作用在态上时，，必须使系统的总电荷减少 $q$。
\end{itemize}
为了协调微观因果性与宏观守恒律，场 $\Psi(x)$ 必须由湮灭部分与产生部分线性叠加而成。我们依次考察这两部分对系统总电荷的贡献，确定粒子的性质。若种类$n$的粒子对于荷算符有本征值$q$:
\begin{subequations}
    \begin{align}
        &[Q,a(\boldsymbol{p},\sigma,n)]=-qa(\boldsymbol{p},\sigma,n)\\
        &[Q,a^\dagger(\boldsymbol{p},\sigma,n)]=+qa^\dagger(\boldsymbol{p},\sigma,n)
    \end{align}
\end{subequations}
对于相互作用$\mathscr{H}\propto \Psi^\dagger_{\ell_1}\Psi^\dagger_{\ell_2}\times\dots\Psi_{\bar{\ell}_1}\Psi_{\bar{\ell}_2}\times\dots$这样的形式，，守恒的要求是：
\begin{equation}
    q_{\ell_1}+q_{\ell_2}+\dots-q_{\bar{\ell}_2}-q_{\bar{\ell}_2}-\dots=0
\end{equation}
为了得到这样的结果，\color{blue}理论是要求自带反粒子；或者说理论本身不带荷\normalcolor，所以将一般形式的场写为：\begin{equation}
    \Psi_\ell(x) =(\frac{1}{2\pi})^\frac{3}{2}\sum_{\sigma}\int d^3p \, \left\{\xi_\ell u_\ell(\boldsymbol{p},\sigma)a^{c\dagger}(\mathbf{p},\sigma) e^{-ipx} + \lambda_\ell v_\ell(\boldsymbol{p},\sigma) a(\mathbf{p},\sigma)e^{ipx} \right\}
\end{equation}
其中
\begin{equation}
    a^{c\dagger}(\mathbf{p},\sigma,n)=a^{\dagger}(\mathbf{p},\sigma,\bar{n})
\end{equation}
$c$是$charge \,conjugate$，若理论本身无荷，，是$a^{c\dagger}(\mathbf{p},\sigma)=a^{\dagger}(\mathbf{p},\sigma)$
